\vspace{-2mm}
\section{Discussion and Conclusion}
\cut{
We introduced an adversarial framework to enforce fairness constraints on graph embeddings. We present two different strategies two enforce fairness: (1) a non-compositional adversary and (2) a compositional adversary. Each type of adversary critiques the representations learned by the main model such that sensitive attributes are not encoded by the main model.
In the compositional approach, we learn a bank of adversarial filters to optionally enforce fairness constraints over a set of possible sensitive attributes.
}

Our work sheds light on how fairness can be enforced in graph representation learning---a setting that is highly relevant to large-scale social recommendation and networking platforms. 
We found that using our proposed compositional adversary allows us to flexibly accomodate unseen combinations of fairness constraints without explicitly training on them. This highlights how fairness could be deployed in a real-word, user-driven setting, where it is necessary to optionally enforce a large number of possible invariance constraints over learned graph representations. 
%We empirically validated our approach on three standard datasets and find that is indeed possible to enforce individual fairness for sets of sensitive attribute, with a performance tradeoff that we illustrate empirically.

In terms of limitations and directions for future work, one important limitation is that we only consider one type of adversarial loss to enforce fairness.
While this adversarial loss is theoretically motivated and known to perform well, there are other recent variations in the literature (e.g., \citet{madras2018learning})---as well as related non-adversarial regularizers (e.g., \citet{zemel2013learning}). Also, while we considered imposing fairness over sets of attributes, we did not explicitly model subgroup-level fairness \cite{kearns2017preventing}.
Extending and testing our framework with these alternatives is a natural direction for future work. 

There are also important questions about how our framework translates to real-world production systems. For instance, in this work we enforced fairness with respect to randomly sampled sets of attributes, but in real-world environments, these sets of attributes would come from user preferences, which may themselves be biased; e.g., it might be more common for female users to request fairness than male users potentially leading to new kinds of demographic inequalities. Understanding how these preference biases could impact our framework is an important direction for future inquiry. 



